\section{Experimental Calibration and Regression}
\label{sec:astrophysical_experiment}

This section tests the SPADE receiver described above using the tabletop configuration and its physical noise parametrization.

\subsection{Regression Calibration Models}
Recall that the average detected fraction $f_1 = n_1/(\eta n)$ incorporates physical cross talk ($\chi = 0.0035$) and stochastic misalignment into the effective noise parameter $\sigma_{\text{eff}}^2 \coloneqq \sigma_M^2 + 4w_0^2\chi$ (Eq.~\eqref{eq:sigma_effective}).

To construct the experimental calibration curves, data were collected by translating the B beam between $-200\,\mu\text{m}$ and $200\,\mu\text{m}$ ($d_a = d/w_0 \in [-0.67, 0.67]$). For each step, we recorded the combined photon counts $n_1 = C^n_{HG01} + C^n_{HG10}$ during $1\text{ s}$ of acquisition time, composed of $N_m = 100$ independent measurements of $10\text{ ms}$ each.

\begin{figure}[h!]
    \centering
    \begin{center}
        \textbf{Calibration Curves for Separation $d_a$ at Four Intensity Ratios $\epsilon$:} \\
        \begin{tabular}{cc}
            $\epsilon = 0.018: \quad n_1^F(d_a) = a_1 + b_1 d_a^2$ & $\epsilon = 0.042: \quad n_1^F(d_a) = a_1 + b_1 d_a^2$ \\
            $\epsilon = 0.069: \quad n_1^F(d_a) = a_1 + b_1 d_a^2$ & $\epsilon = 0.096: \quad n_1^F(d_a) = a_1 + b_1 d_a^2$
        \end{tabular}
    \end{center}
    \caption{Quadratic regression curves of the first-order mode counts $n_1(d_a)$ versus dimensionless separation $d_a$ for the four experimental intensity ratios $\epsilon \in \{0.018, 0.042, 0.069, 0.096\}$. The fitting parameters $a_1$ and $b_1$ account for the physical cross talk ($\chi = 0.0035$) and dark count rates.}
    \label{fig:calibration_da_link}
\end{figure}

The fitting parameter $a_1$ corresponds to the baseline count rate from the effective noise parameter $\sigma_{\text{eff}}^2$, while $b_1$ represents the parameter-dependent signal scaling:
\begin{align}
    a_1 & \propto \eta \sigma_{\text{eff}}^2 = \eta (\sigma_M^2 + 4w_0^2\chi), \\
    b_1 & \propto \eta \epsilon w_0^2.
\end{align}
The resulting experimental uncertainties $\Delta d_a$ obtained via these quadratic calibrations show that the SPADE tabletop setup beats the direct-imaging theoretical limits for separations down to $d_a \approx 0.15$.

For the relative intensity ratio $\epsilon$, linear calibrations were performed at the four fixed separations $d_a \in \{0.2, 0.27, 0.47, 0.67\}$:

\begin{figure}[h!]
    \centering
    \begin{center}
        \textbf{Calibration Curves for Intensity $\epsilon$ at Four Separations $d_a$:} \\
        \begin{tabular}{cc}
            $d_a = 0.2: \quad n_1^F(\epsilon) = \alpha_1 + \beta_1 \epsilon$  & $d_a = 0.27: \quad n_1^F(\epsilon) = \alpha_1 + \beta_1 \epsilon$ \\
            $d_a = 0.47: \quad n_1^F(\epsilon) = \alpha_1 + \beta_1 \epsilon$ & $d_a = 0.67: \quad n_1^F(\epsilon) = \alpha_1 + \beta_1 \epsilon$
        \end{tabular}
    \end{center}
    \caption{Linear regression calibration curves of the first-order mode counts $n_1(\epsilon)$ versus relative intensity $\epsilon$ for the four fixed dimensionless separations $d_a \in \{0.2, 0.27, 0.47, 0.67\}$.}
    \label{fig:calibration_eps_link}
\end{figure}

The fitting parameter $\alpha_1$ accounts for the background counts from spatial crosstalk, while $\beta_1$ scales linearly with the separation $d_a^2$. The resulting experimental uncertainties $\Delta\epsilon$ match the theoretical stochastic error models (Eqs.~\eqref{eq:propagation_uncertainty_d}--\eqref{eq:std_dev_fraction}) within a factor of two.

\section{Experimental Channel Model and Overdispersed Counts}
\label{sec:practical_channel_framework}

To analyze the raw experimental dataset, we model the apparatus as a cascade of three physical and statistical channels. The model keeps the partially degenerate experimental parameters visible, fits the photon survival rate $\eta$ using maximum-likelihood estimation (MLE), and assesses the resulting count distributions using parametric bootstrapping.

\subsection{The Three-Stage Sequential Channel Cascade}

The tabletop experiment maps the known input controls, the dimensionless separation $d_a = d/w_0 \in [-0.67, 0.67]$ and relative intensity ratio $\epsilon \in [0.004, 0.096]$, to the raw detected counts $n_1$ on the $HG_{01} + HG_{10}$ detectors through the following sequential channel cascade:

\begin{figure}[H]
    \centering
    \begin{center}
        \textbf{The Laboratory Channel Pipeline:} \\
        \begin{tabular}{ccccc}
            \colorbox{lightgray}{\textbf{Input Controls}} & $\xrightarrow{\quad\mathcal{C}_{\text{state}}\quad}$ & \colorbox{lightgray}{\textbf{Quantum State}} & $\xrightarrow{\quad\mathcal{C}_{\text{readout}}\quad}$ & \colorbox{lightgray}{\textbf{Physical Mode Sorting}} \\
            $(d_a, \epsilon)$                             &                                                      & $\rho_d(d_a, \epsilon)$                      &                                                        & Pre-Detector Probabilities $p_0, p_1$                \\
        \end{tabular}
        \\[0.3cm]
        \begin{tabular}{ccc}
            $\xrightarrow{\quad\mathcal{C}_{\text{count}}\quad}$ & \colorbox{lightgray}{\textbf{Overdispersed Count Noise}} & $\longrightarrow \quad \mathbf{n_1 \text{ counts (Raw Data)}}$ \\
                                                                 & (Fano Factor $F \approx 2.6$)                            &                                                                \\
        \end{tabular}
    \end{center}
    \caption{Schematic representation of the sequential physical and statistical channels modeling the tabletop SPADE imaging apparatus.}
    \label{fig:sequential_channel_pipeline}
\end{figure}

\begin{enumerate}
    \item \textbf{Quantum state channel ($\mathcal{C}_{\text{state}}$).} This channel maps the classical coordinates $(d_a, \epsilon)$ to the single-photon density matrix representing the mixed astronomical light field (Eq.~\eqref{eq:optical_density_matrix_chap}):
          \begin{equation}
              \rho_d(d_a, \epsilon) = (1 - \eta) |0\rangle\langle 0| + \eta (1 - \epsilon) |\psi_0\rangle\langle\psi_0| + \eta \epsilon |\psi_d\rangle\langle\psi_d|.
          \end{equation}
          Here $\eta$ is the joint probability of single-photon generation, transmission, and detection.

    \item \textbf{Optical readout channel ($\mathcal{C}_{\text{readout}}$).} This channel models geometric alignment, aperture losses, and mode-sorting crosstalk through
          \begin{equation}
              \begin{pmatrix} p_0^{\text{real}} \\ p_1^{\text{real}} \end{pmatrix} =
              \begin{pmatrix} 1 - \chi & l_0 \\ \chi & 1 - l_0 \end{pmatrix}
              \begin{pmatrix} p_0 \\ p_1 \end{pmatrix} +
              \begin{pmatrix} 0 \\ l_{\text{tail}} p_{\text{tail}} \end{pmatrix}.
          \end{equation}
          We use $l_0 \approx 0.002$ for leakage of $HG_{00}$ into the first-order output ports, $\chi = 0.0035$ for demultiplexer crosstalk, and $l_{\text{tail}} \approx 0.002$ for higher-order leakage.

    \item \textbf{Count-noise channel ($\mathcal{C}_{\text{count}}$).} This final classical channel maps $p_1^{\text{real}}$ to the recorded counts $n_1$. The data have a median Fano factor
          \begin{equation}
              F \coloneqq \frac{\operatorname{Var}(n_1)}{\mathbb{E}[n_1]} \approx 2.50,
          \end{equation}
          with a $5^{\text{th}}$--$95^{\text{th}}$ percentile range of $1.85$--$3.39$. We therefore use the Fano-corrected variance model
          \begin{equation}
              \operatorname{Var}(n_1) = F \cdot \mathbb{E}[n_1], \qquad F = 2.6.
              \label{eq:fano_overdispersion}
          \end{equation}
\end{enumerate}

\subsection{Maximum-Likelihood Estimation and Parametric Bootstrap}

Fixing the incident flux $N_{\text{in}} = 40{,}000$ photons, crosstalk $\chi = 0.0035$, and leakage $l_0 = 0.002$, we use MLE to fit the single-photon survival rate $\eta$.

\begin{figure}[H]
    \centering
    \begin{center}
        \textbf{Calibration Performance under Noisy Channels:}
        \begin{tabular}{cccc}
            \hline\hline
            \textbf{Noise Model}               & \textbf{MLE Survival Rate $\eta$} & \textbf{$99.9\%$ Profile Likelihood} & \textbf{Bootstrap $p$-value}                \\
            \hline
            Poissonian ($F = 1$)               & $0.11258$                         & ---                                  & $0.00498 \quad \text{(Rejected)}$           \\
            \textbf{Overdispersed ($F = 2.6$)} & $\mathbf{0.10480}$                & $\mathbf{[0.10408, \, 0.10553]}$     & $\mathbf{0.80600 \quad \text{(Validated)}}$ \\
            \hline\hline
        \end{tabular}
    \end{center}
    \caption{Comparison of Poissonian and overdispersed count-noise models fitted to the raw SPADE dataset.}
    \label{fig:mle_bootstrap_table}
\end{figure}

The Poisson model is rejected by the bootstrap test ($p=0.00498$). Incorporating the Fano channel yields $\eta_{\text{MLE}} = 0.10480$, with a $99.9\%$ profile-likelihood interval of $[0.10408, 0.10553]$ and bootstrap $p=0.8060$.

\subsection{QFI on the Pre-Detector State}

The overdispersed count noise is a detector-level statistical effect, so the QFI is evaluated on the pre-detector physical state $\rho_{\text{pre}}(d_a, \epsilon)$ after optical readout but before count detection:
\begin{equation}
    \mathcal{I}_{\text{pre}}(d) \coloneqq \operatorname{Tr}\left[ \rho_{\text{pre}}(d) L_d^2 \right].
    \label{eq:qfi_pre_detector}
\end{equation}
Here $\rho_{\text{pre}}(d)$ incorporates the fitted attenuation $\eta_{\text{MLE}}$, crosstalk $\chi$, and modal leakage $l_0$. Using the vectorized formulation of Eq.~\eqref{eq:safranek_vectorized}, the model gives
\begin{equation}
    \lim_{d \to 0} \mathcal{I}_{\text{pre}}(d) = \frac{\eta_{\text{MLE}} \epsilon}{w_0^2} \left( 1 - \frac{\chi}{\epsilon + \sigma^2_{\text{eff}}/4w_0^2} \right).
\end{equation}
This connects the measured crosstalk and leakage parameters to the experimentally accessible precision of the SPADE receiver.
